% ch54_triangulated_t_structures.tex — Volume IV Chapter 18
% Retrofitted with full Vol I-III pedagogy: cycle stages, etymology, dialogs.

\chapter{Triangulated Categories and t-Structures}

\begin{overviewbox}
The homotopy category $\mathrm{ho}(\mathcal{C})$ of a stable
$\infty$-category $\mathcal{C}$ comes equipped with an extra layer of
structure: it is a \term{triangulated category}, in the sense of Verdier.
The triangulated structure remembers what the $\infty$-categorical fibre/cofibre
sequences look like at the $1$-categorical level: distinguished triangles
$a \to b \to c \to \Sigma a$ replace the original homotopical data.

This chapter has two main themes.

First, the construction $\mathrm{ho}(\mathrm{stable}) \to \mathrm{triangulated}$.
We define triangulated categories axiomatically (Verdier's axioms TR1–TR4),
verify that $\mathrm{ho}(\mathcal{C})$ satisfies them, and explore the
classical example: the homotopy category of chain complexes
$\mathbf{K}(\mathcal{A})$ together with its triangulated structure.

Second, \term{t-structures}: a pair of full subcategories $(\mathcal{C}_{\geq 0},
\mathcal{C}_{\leq 0})$ of a stable $\infty$-category (or its triangulated
homotopy category) satisfying conditions that ensure
\begin{itemize}
  \item Each object decomposes into ``connective and coconnective parts.''
  \item The \term{heart} $\mathcal{C}^\heartsuit := \mathcal{C}_{\geq 0}
        \cap \mathcal{C}_{\leq 0}$ is an abelian category.
\end{itemize}

The motivating example: the standard t-structure on
$\mathcal{D}(\mathcal{A})$ has heart $\mathcal{A}$ itself, and the
truncation functors recover the standard cohomology
$H^n : \mathcal{D}(\mathcal{A}) \to \mathcal{A}$. Other t-structures
(perverse, étale, motivic) produce different ``hearts'' suitable for
specific applications.

Together, stable $\infty$-categories and t-structures supply the
$\infty$-categorical home for homological algebra and its applications.
\end{overviewbox}


% ═══════════════════════════════════════════════════════════════════════════
\section{Triangulated Categories}
\label{sec:triangulated}
% ═══════════════════════════════════════════════════════════════════════════

\subsection{Formalizing}

\begin{nameorigin}
\term{Triangulated category} was introduced by \emph{Jean-Louis Verdier}
(1935--1989) in his 1967 PhD thesis (advisor: Grothendieck). The name
comes from \emph{distinguished triangles} — sequences $a \to b \to c
\to \Sigma a$ that play the role in triangulated categories that exact
sequences play in abelian categories. The triangle is a tightening of
the long exact sequence $\cdots \to a \to b \to c \to \Sigma a \to \cdots$
into a single rotating triple. Verdier's axiomatization captures the
structure of derived categories of abelian categories abstractly; the
$\infty$-categorical lift via stable $\infty$-cats was Lurie (HA 2017).
\end{nameorigin}

\begin{nameorigin}
The \term{octahedron axiom} (TR4) is the most distinctive — and most
delicate — axiom of triangulated categories. The name comes from the
geometric shape of the diagram: given two composable morphisms $a \to b \to c$,
the resulting four distinguished triangles fit into an octahedron with
eight triangular faces. The axiom expresses a compatibility between the
three composable cones. Verdier's original 1967 formulation; many
authors have rewritten the axiom over the years (e.g., Beilinson-Bernstein-Deligne
introduced a cleaner $9$-diagram version).
\end{nameorigin}

\begin{definition}[Triangulated category]
\label{def:triangulated}
A \term{triangulated category} is an additive category $\mathcal{T}$
equipped with:
\begin{itemize}
  \item A \term{shift} or \term{suspension} autoequivalence $\Sigma :
        \mathcal{T} \xrightarrow{\sim} \mathcal{T}$.
  \item A class of \term{distinguished triangles} of the form
        $a \to b \to c \to \Sigma a$,
\end{itemize}
satisfying Verdier's axioms \term{TR1}--\term{TR4}:
\begin{itemize}
  \item[\textbf{TR1}] (Identity) For every $a$, the triangle $a \xrightarrow{\mathrm{id}}
        a \to 0 \to \Sigma a$ is distinguished. Distinguished triangles
        are closed under isomorphism.
  \item[\textbf{TR2}] (Rotation) $a \to b \to c \to \Sigma a$ is distinguished
        iff $b \to c \to \Sigma a \xrightarrow{-\Sigma f} \Sigma b$ is.
  \item[\textbf{TR3}] (Morphisms of triangles) Given two distinguished
        triangles and a morphism between their first two terms, there
        is a (not necessarily unique) extension to a morphism between
        the full triangles.
  \item[\textbf{TR4}] (Octahedron axiom) Given two composable morphisms
        $a \to b \to c$, the resulting distinguished triangles fit into
        an ``octahedron'' of compatible distinguished triangles.
\end{itemize}
\end{definition}

\begin{howtosay}
Distinguished triangle \sayit{``distinguished triangle'' or ``exact triangle''}.
The notation $a \to b \to c \to \Sigma a$ implicitly carries the data of
all three connecting morphisms.
\end{howtosay}

\subsection{Reorganizing: where triangulated comes from}

\begin{theorem}[$\mathrm{ho}$ of stable is triangulated]
\label{thm:ho-stable-triangulated}
For any stable $\infty$-category $\mathcal{C}$, the homotopy $1$-category
$\mathrm{ho}(\mathcal{C})$ is naturally triangulated. The shift is induced
by $\Sigma$; the distinguished triangles are the images in $\mathrm{ho}$
of the (co)fibre sequences in $\mathcal{C}$.
\end{theorem}

\begin{proof}[Proof strategy]
Verifying TR1–TR4 amounts to constructing simplices in $\mathcal{C}$:
\begin{itemize}
  \item TR1: the trivial fibre sequence $a \xrightarrow{\mathrm{id}} a \to 0$ provides the
        identity triangle.
  \item TR2 (rotation) is the iterated cofibre construction.
  \item TR3 comes from inner-horn filling in $\mathcal{C}$.
  \item TR4 (octahedron) corresponds to a specific $3$-simplex construction
        in $\mathcal{C}$ encoding compatibility of two cofibres.
\end{itemize}
All steps use the stability axiom (pullbacks $=$ pushouts) at the
$\infty$-categorical level.
\qed
\end{proof}

\begin{example_thm}[$\mathbf{K}(\mathcal{A})$ and $\mathcal{D}(\mathcal{A})$]
\label{ex:K-D-triangulated}
For an abelian category $\mathcal{A}$:
\begin{itemize}
  \item $\mathbf{K}(\mathcal{A})$ (chain complexes modulo chain homotopy)
        is triangulated: shift by one, mapping cone gives distinguished
        triangles.
  \item $\mathcal{D}(\mathcal{A})$ (the derived category) is triangulated:
        the same structure, descending through localization at
        quasi-isomorphisms.
\end{itemize}
$\mathcal{D}(\mathcal{A}) = \mathrm{ho}$ of the stable $\infty$-category
$\mathcal{D}^\infty(\mathcal{A})$.
\end{example_thm}


% ═══════════════════════════════════════════════════════════════════════════
\section{t-Structures}
\label{sec:t-structures}
% ═══════════════════════════════════════════════════════════════════════════

\subsection{Formalizing}

\begin{nameorigin}
\term{t-structure} was introduced by \emph{Beilinson, Bernstein, and
Deligne} (\emph{Faisceaux Pervers}, 1982) in their development of
perverse sheaves. The letter ``t'' originally stood for \emph{topology}
— BBD's t-structures generalize the Postnikov truncation of topological
spaces, where the bounded-above and bounded-below subcategories
correspond to ``everything above'' and ``everything below'' a given
degree. The \emph{heart} of a t-structure (the intersection of the two
parts) is an abelian category — a remarkable recovery of abelian
structure from triangulated data. The terminology has stuck despite
the more modern view that t-structures generalize \emph{many} kinds of
truncations (Postnikov, étale, perverse, etc.).
\end{nameorigin}

\begin{definition}[t-structure]
\label{def:t-structure}
A \term{t-structure} on a stable $\infty$-category $\mathcal{C}$ is a
pair of full subcategories $(\mathcal{C}_{\geq 0}, \mathcal{C}_{\leq 0})$
of $\mathcal{C}$ satisfying:
\begin{enumerate}
  \item $\Sigma \mathcal{C}_{\geq 0} \subseteq \mathcal{C}_{\geq 0}$ and
        $\Omega \mathcal{C}_{\leq 0} \subseteq \mathcal{C}_{\leq 0}$.
  \item For $x \in \mathcal{C}_{\geq 0}$ and $y \in \Sigma^{-1} \mathcal{C}_{\leq 0}$:
        $\mathrm{Map}_\mathcal{C}(x, y) \simeq 0$ (mapping space is contractible).
  \item Every $z \in \mathcal{C}$ fits into a fibre sequence
        $\tau_{\geq 0} z \to z \to \tau_{\leq -1} z$
        with $\tau_{\geq 0} z \in \mathcal{C}_{\geq 0}$ and $\tau_{\leq -1} z
        \in \mathcal{C}_{\leq -1} = \Sigma^{-1} \mathcal{C}_{\leq 0}$.
\end{enumerate}
We define $\mathcal{C}_{\geq n} = \Sigma^n \mathcal{C}_{\geq 0}$ and
$\mathcal{C}_{\leq n} = \Sigma^n \mathcal{C}_{\leq 0}$.
\end{definition}

\begin{howtosay}
$\mathcal{C}_{\geq 0}$ \sayit{``the connective part of $\mathcal{C}$''} or
\sayit{``the $\geq 0$-truncated part''}. $\mathcal{C}_{\leq 0}$:
\sayit{``coconnective''}. $\tau_{\geq 0}, \tau_{\leq 0}$ are
\term{truncation functors}.
\end{howtosay}

\begin{nameorigin}
The \term{heart} $\mathcal{C}^\heartsuit$ of a t-structure is the
abelian category at the ``center'' of the triangulated structure — the
intersection of the connective and coconnective parts. The terminology
is BBD (1982): the heart is the ``heart'' of the structure, the central
part from which everything else can be reconstructed via cohomology
functors. The symbol $\heartsuit$ (heart suit) is universal across the
literature.
\end{nameorigin}

\begin{theorem}[Heart of a t-structure]
\label{thm:heart}
The intersection $\mathcal{C}^\heartsuit := \mathcal{C}_{\geq 0} \cap
\mathcal{C}_{\leq 0}$ is an \emph{abelian category}. It is closed under
finite limits and finite colimits in $\mathcal{C}$, but not closed under
suspension (which takes $\mathcal{C}^\heartsuit$ outside itself).
\end{theorem}

\begin{example_thm}[Standard t-structure on $\mathcal{D}(\mathcal{A})$]
\label{ex:standard-t-D}
For an abelian category $\mathcal{A}$, the derived $\infty$-category
$\mathcal{D}(\mathcal{A})$ has a canonical t-structure:
\begin{itemize}
  \item $\mathcal{D}(\mathcal{A})_{\geq 0}$: chain complexes whose
        $H^n = 0$ for $n < 0$.
  \item $\mathcal{D}(\mathcal{A})_{\leq 0}$: chain complexes whose
        $H^n = 0$ for $n > 0$.
  \item Heart: $\mathcal{D}(\mathcal{A})^\heartsuit \simeq \mathcal{A}$ via
        the cohomology functor $H^0$.
\end{itemize}
\end{example_thm}

\subsection{Reorganizing: cohomology functors}

\begin{definition}[Cohomology functor]
\label{def:cohomology}
The \term{$n$-th cohomology functor} of a t-structure is
\begin{equation*}
  H^n : \mathcal{C} \to \mathcal{C}^\heartsuit, \quad
  H^n(x) = \tau_{\leq 0} \tau_{\geq 0} (\Omega^n x).
\end{equation*}
$H^n$ is exact in the sense that it sends fibre/cofibre sequences in
$\mathcal{C}$ to long exact sequences in $\mathcal{C}^\heartsuit$.
\end{definition}

\begin{theorem}[Long exact sequence in cohomology]
\label{thm:LES-cohomology}
A fibre sequence $a \to b \to c$ in $\mathcal{C}$ induces a long exact
sequence in the heart $\mathcal{C}^\heartsuit$:
\begin{equation*}
  \cdots \to H^{n-1}(c) \to H^n(a) \to H^n(b) \to H^n(c) \to H^{n+1}(a) \to \cdots.
\end{equation*}
\end{theorem}


% ═══════════════════════════════════════════════════════════════════════════
\section{Other t-Structures}
\label{sec:other-t}
% ═══════════════════════════════════════════════════════════════════════════

\subsection{Formally: variants}

\begin{example_thm}[Perverse t-structure]
\label{ex:perverse}
On the derived category of constructible sheaves on a topological space,
the \term{perverse t-structure} (Beilinson–Bernstein–Deligne) has
distinct connective and coconnective parts from the standard one. Its
heart consists of \term{perverse sheaves}, an abelian category
playing a central role in geometric representation theory and the
geometric Langlands program.
\end{example_thm}

\begin{example_thm}[Étale t-structure]
\label{ex:etale}
On the derived $\infty$-category of étale sheaves on a scheme, a different
t-structure (the standard one for étale cohomology) has heart the
category of étale sheaves of abelian groups.
\end{example_thm}

\begin{example_thm}[Postnikov t-structure on spectra]
\label{ex:postnikov-spectra}
On $\mathrm{Sp}$, the \term{Postnikov t-structure} has connective spectra
($\pi_n = 0$ for $n < 0$) and coconnective spectra ($\pi_n = 0$ for $n > 0$).
Heart: $\mathbf{Ab}$, the category of abelian groups, via Eilenberg–MacLane
spectra.
\end{example_thm}


% ═══════════════════════════════════════════════════════════════════════════
\section{Exercises}
\label{sec:triangulated-exercises}
% ═══════════════════════════════════════════════════════════════════════════

\begin{exercisesbox}
\begin{qa}
\qaitem{Define a triangulated category.}{An additive category with a shift autoequivalence $\Sigma$ and a class of distinguished triangles, satisfying Verdier's TR1--TR4.}
\qaitem{State TR2 (rotation).}{$a \to b \to c \to \Sigma a$ distinguished iff $b \to c \to \Sigma a \xrightarrow{-\Sigma f} \Sigma b$ distinguished — rotating triangles preserves the distinguished property up to a sign.}
\qaitem{State TR4 (octahedron).}{Given $a \to b \to c$ as composable morphisms, the resulting four distinguished triangles fit into an octahedron of compatible morphisms.}
\qaitem{State the theorem connecting stable $\infty$-cats and triangulated categories.}{For any stable $\infty$-cat $\mathcal{C}$, $\mathrm{ho}(\mathcal{C})$ is canonically triangulated. Distinguished triangles are images of fibre/cofibre sequences in $\mathcal{C}$.}
\qaitem{Why does $\mathrm{ho}(\mathcal{C})$ acquire triangulated structure?}{Because stability of $\mathcal{C}$ provides the symmetry between fibre and cofibre sequences (TR2); inner-horn filling provides morphism extensions (TR3); and a specific $3$-simplex realizes octahedron (TR4).}
\qaitem{What's $\mathbf{K}(\mathcal{A})$?}{The category of chain complexes modulo chain homotopy — the most classical example of a triangulated category. Distinguished triangles arise from mapping cones.}
\qaitem{Define a t-structure.}{A pair $(\mathcal{C}_{\geq 0}, \mathcal{C}_{\leq 0})$ of subcategories satisfying closure under suspension/loop, mapping-space orthogonality, and a truncation triangle for every object.}
\qaitem{Define the heart.}{$\mathcal{C}^\heartsuit = \mathcal{C}_{\geq 0} \cap \mathcal{C}_{\leq 0}$ — an abelian category.}
\qaitem{What's the standard t-structure on $\mathcal{D}(\mathcal{A})$?}{$\mathcal{D}(\mathcal{A})_{\geq 0}$ = complexes with $H^n = 0$ for $n < 0$; $\mathcal{D}(\mathcal{A})_{\leq 0}$ = complexes with $H^n = 0$ for $n > 0$. Heart: $\mathcal{A}$.}
\qaitem{What's the cohomology functor?}{$H^n(x) = \tau_{\leq 0} \tau_{\geq 0} (\Omega^n x)$. Takes a stable object to its $n$-th cohomology in the heart.}
\qaitem{State the long exact sequence in cohomology.}{For a fibre sequence $a \to b \to c$: a LES $\cdots \to H^{n-1}(c) \to H^n(a) \to H^n(b) \to H^n(c) \to \cdots$ in the heart.}
\qaitem{What's the Postnikov t-structure on $\mathrm{Sp}$?}{Connective spectra (positive $\pi_n$ only) vs coconnective (negative $\pi_n$ only). Heart: $\mathbf{Ab}$ via Eilenberg–MacLane spectra.}
\qaitem{What's a perverse t-structure?}{A non-standard t-structure on derived categories of constructible sheaves whose heart is perverse sheaves. Central to geometric representation theory.}
\qaitem{Why don't all stable $\infty$-categories have natural t-structures?}{Because t-structures encode a specific choice of ``positive/negative'' decomposition that not every stable category admits canonically. Spectra have one; abstract triangulated categories may not.}
\qaitem{Is the heart of a t-structure abelian?}{Yes (Theorem~\ref{thm:heart}). The conditions on the t-structure guarantee exactness of finite limits and colimits within the heart.}
\qaitem{Are stable $\infty$-categories ``more structured'' than triangulated ones?}{Yes — stable $\infty$-cats retain higher-coherence data lost in passing to $\mathrm{ho}$. Triangulated categories are the $1$-truncation; some constructions (e.g., octahedron) become non-canonical at the $1$-level.}
\qaitem{Why is the octahedron axiom delicate?}{Because the choice of morphisms in TR4 is non-unique at the triangulated level; only the $\infty$-categorical lift makes it functorial. This is a source of incompleteness in classical triangulated theory.}
\qaitem{What's the relation to derived functors of Chapter~42?}{Derived functors are exact functors between stable $\infty$-cats (or triangulated cats). Tor and Ext (Ch~42) compute cohomology in standard t-structures.}
\qaitem{What does Riehl--Verity / Lurie use stable $\infty$-cats for primarily?}{Stable homotopy theory, derived algebraic geometry, modular representation theory, and the modern theory of categorified algebra (e.g., 2-categorical structures arising from stable diagram homotopy).}
\qaitem{What's Verdier's theorem on triangulated functors?}{Functors between triangulated categories preserving shifts and distinguished triangles correspond, $\infty$-categorically, to exact functors. So $\mathrm{ho}$ of an exact functor of stable $\infty$-cats is a triangulated functor.}
\qaitem{Is every triangulated category $\mathrm{ho}$ of a stable $\infty$-cat?}{Not obviously: there exist triangulated categories with no known stable $\infty$-categorical lift (``ghost-classified'' triangulated categories). It's a famous open problem in higher algebra.}
\qaitem{What's a t-exact functor?}{A functor between stable $\infty$-cats with t-structures that preserves the connective and coconnective parts: $F(\mathcal{C}_{\geq 0}) \subseteq \mathcal{D}_{\geq 0}$ and similarly for $\leq 0$.}
\qaitem{What's the next chapter?}{Ch~55: $\infty$-Operads. Moving from stable to higher-algebraic structures: the operadic language for organizing multilinear and commutative-up-to-coherent operations.}
\end{qa}
\end{exercisesbox}


% ═══════════════════════════════════════════════════════════════════════════
\section*{Chapter Summary}
\addcontentsline{toc}{section}{Chapter Summary}
% ═══════════════════════════════════════════════════════════════════════════

\begin{summarybox}
\textbf{Triangulated category.}\quad Additive category with shift
$\Sigma$ and distinguished triangles, satisfying Verdier's TR1--TR4
(identity, rotation, morphisms-extend, octahedron).

\textbf{$\mathrm{ho}(\mathrm{stable})$ is triangulated.}\quad Every
stable $\infty$-cat's homotopy category is triangulated, with
distinguished triangles from fibre/cofibre sequences.

\textbf{Classical examples.}\quad $\mathbf{K}(\mathcal{A})$, $\mathcal{D}(\mathcal{A})$:
chain complexes modulo chain homotopy, and the derived category.

\textbf{t-structure.}\quad Pair $(\mathcal{C}_{\geq 0}, \mathcal{C}_{\leq 0})$
with closure, orthogonality, truncation. Encodes a ``positive/negative''
decomposition. Heart $\mathcal{C}^\heartsuit$ is abelian.

\textbf{Cohomology.}\quad $H^n : \mathcal{C} \to \mathcal{C}^\heartsuit$
takes stable objects to abelian cohomology. Sends fibre sequences to
long exact sequences.

\textbf{Variants.}\quad Standard t-structure on $\mathcal{D}(\mathcal{A})$
with heart $\mathcal{A}$; Postnikov on $\mathrm{Sp}$ with heart $\mathbf{Ab}$;
perverse for constructible sheaves; étale for étale sheaves.
\end{summarybox}


% ═══════════════════════════════════════════════════════════════════════════
\section*{Formal Restatement}
\addcontentsline{toc}{section}{Formal Restatement}
% ═══════════════════════════════════════════════════════════════════════════

\begin{formalrestatement}
\textbf{Triangulated category.}\quad $(\mathcal{T}, \Sigma, \text{distinguished triangles})$
satisfying TR1--TR4.

\medskip
\textbf{$\mathrm{ho}$-stable.}\quad $\mathcal{C}$ stable $\Rightarrow
\mathrm{ho}(\mathcal{C})$ canonically triangulated.

\medskip
\textbf{t-structure.}\quad $(\mathcal{C}_{\geq 0}, \mathcal{C}_{\leq 0})$ on a stable
$\mathcal{C}$:
\begin{itemize}
  \item $\Sigma \mathcal{C}_{\geq 0} \subseteq \mathcal{C}_{\geq 0}$, $\Omega \mathcal{C}_{\leq 0} \subseteq \mathcal{C}_{\leq 0}$.
  \item $x \in \mathcal{C}_{\geq 0}, y \in \mathcal{C}_{\leq -1} \Rightarrow \mathrm{Map}(x, y) = 0$.
  \item $\forall z, \exists$ fibre seq $\tau_{\geq 0} z \to z \to \tau_{\leq -1} z$ with $\tau_{\geq 0} z \in \mathcal{C}_{\geq 0}$, $\tau_{\leq -1} z \in \mathcal{C}_{\leq -1}$.
\end{itemize}

\medskip
\textbf{Heart.}\quad $\mathcal{C}^\heartsuit = \mathcal{C}_{\geq 0} \cap \mathcal{C}_{\leq 0}$:
abelian.

\medskip
\textbf{Cohomology.}\quad $H^n = \tau_{\leq 0} \tau_{\geq 0} \Omega^n :
\mathcal{C} \to \mathcal{C}^\heartsuit$. LES from fibre sequences.

\medskip
\textbf{Examples.}\quad
$\mathcal{D}(\mathcal{A})^\heartsuit = \mathcal{A}$;
$\mathrm{Sp}^\heartsuit = \mathbf{Ab}$ (Postnikov);
perverse and étale variants.
\end{formalrestatement}
